\documentclass[tikz,border=8pt]{standalone}
\usepackage{tikz}
\usepackage{amsmath}
\usepackage{pgfplots}
\pgfplotsset{compat=1.18}
\usepackage{ctex}
\usetikzlibrary{calc, positioning, arrows.meta, decorations.pathreplacing}

\begin{document}
\begin{tikzpicture}

%% ===== 左图：平行四边形法则 =====
\begin{scope}[xshift=0cm]

  % 坐标系背景网格
  \draw[gray!25, thin] (-0.3,-0.3) grid (4.3,3.3);
  \draw[->,gray!60, thin] (-0.3,0) -- (4.5,0) node[right,font=\scriptsize] {$x$};
  \draw[->,gray!60, thin] (0,-0.3) -- (0,3.5) node[above,font=\scriptsize] {$y$};

  % u = (3,1), v = (1,2)
  % 平行四边形四个顶点: O=(0,0), A=u=(3,1), B=v=(1,2), C=u+v=(4,3)
  \fill[blue!10] (0,0) -- (3,1) -- (4,3) -- (1,2) -- cycle;
  \draw[blue!40, dashed, thin] (3,1) -- (4,3);
  \draw[blue!40, dashed, thin] (1,2) -- (4,3);

  % 向量 u
  \draw[->,thick,blue!80!black,>=Stealth] (0,0) -- (3,1)
    node[below right,font=\small] {$\mathbf{u}=(3,1)$};
  % 向量 v
  \draw[->,thick,teal,>=Stealth] (0,0) -- (1,2)
    node[left,font=\small] {$\mathbf{v}=(1,2)$};
  % 对角线 u+v
  \draw[->,thick,red!80!black,>=Stealth] (0,0) -- (4,3)
    node[above right,font=\small] {$\mathbf{u}+\mathbf{v}=(4,3)$};

  % 原点标注
  \fill (0,0) circle (1.5pt) node[below left,font=\scriptsize] {$O$};

  % 子图标题
  \node[font=\small\bfseries] at (2.0,-0.9) {平行四边形法则};

\end{scope}

%% ===== 右图：三角形法则 =====
\begin{scope}[xshift=6.2cm]

  % 背景网格
  \draw[gray!25, thin] (-0.3,-0.3) grid (4.3,3.3);
  \draw[->,gray!60, thin] (-0.3,0) -- (4.5,0) node[right,font=\scriptsize] {$x$};
  \draw[->,gray!60, thin] (0,-0.3) -- (0,3.5) node[above,font=\scriptsize] {$y$};

  % u 从 O 出发
  \draw[->,thick,blue!80!black,>=Stealth] (0,0) -- (3,1)
    node[below right,font=\small] {$\mathbf{u}=(3,1)$};
  % v 从 u 的终点出发
  \draw[->,thick,teal,>=Stealth] (3,1) -- (4,3)
    node[right,font=\small] {$\mathbf{v}=(1,2)$};
  % u+v 闭合边
  \draw[->,thick,red!80!black,>=Stealth] (0,0) -- (4,3)
    node[above left,font=\small] {$\mathbf{u}+\mathbf{v}=(4,3)$};

  % 三角形填色
  \fill[red!8] (0,0) -- (3,1) -- (4,3) -- cycle;

  % 原点
  \fill (0,0) circle (1.5pt) node[below left,font=\scriptsize] {$O$};
  \fill (3,1) circle (1.5pt) node[below right,font=\scriptsize] {$(3,1)$};
  \fill (4,3) circle (1.5pt) node[above right,font=\scriptsize] {$(4,3)$};

  % 子图标题
  \node[font=\small\bfseries] at (2.0,-0.9) {三角形法则};

\end{scope}

%% 整体标题
\node[font=\large\bfseries] at (6.1, 4.2) {向量加法的两种几何方法};

\end{tikzpicture}
\end{document}
